\documentclass[12pt]{article}
\usepackage[utf8]{inputenc}
\usepackage[T1]{fontenc}
\usepackage{hyperref}
\title{Paipai Language Resources for K-12 Curriculum Development}
\author{Compiled by Claude Sonnet 5.5 (claude-sonnet-5-5) for Andr\'as Kornai}
\date{}
\begin{document}
\maketitle

\section{Introduction}
Paipai (Pai Pai, Akwa'ala; ISO 639-3 ppi) is a critically endangered Cochim\'i-Yuman language of Baja California, Mexico. Published speaker counts are very small (about 25 in Santa Catarina in 2010 and 49 in 2015, according to the Endangered Languages Project). This report is short because few resources surfaced in search.

\section{Primary Educational Resources}
No textbook, primer or school curriculum was found in search results. The Endangered Languages Project page lists an entry titled ``An alphabet of the Paipai language'', but its content was not examined.
\begin{itemize}
\item Endangered Languages Project, Paipai language page: \url{https://endangeredlanguages.com/lang/1538}
\end{itemize}

\section{Linguistic and Reference Materials}
\begin{itemize}
\item Ivette S. Gonz\'alez, ``Language Snapshot: Paipai'', \textit{Language Documentation and Description} vol.\ 17 (2020), an overview of the language and its documentation: \url{https://www.lddjournal.org/article/1210/galley/2454/download}
\item Endangered Languages Project record (alternate entry): \url{https://endangeredlanguages.com/node/112156}
\end{itemize}
No dictionary was found.

\section{Religious and Community Resources}
None found in search results.

\section{Digital and Audio Resources}
The Endangered Languages Project pages above are the only digital entry points found; the snapshot article mentions some audio-visual materials but they were not located directly.

\section{Curriculum Development Context}
With very few, mostly elderly speakers and no found school materials, any curriculum would need to begin with documentation and community-led materials development. Mexico's INALI catalogue was mentioned in search summaries as a further starting point but was not itself retrieved.
\end{document}